% !TEX program = xelatex
% !TEX encoding = UTF-8
% !BIB program = bibtex
% !TEX spellcheck = en_US
\documentclass[12pt]{article}
\usepackage{url}
\usepackage{amsmath}
\usepackage{graphicx}
\usepackage{subfigure}
\graphicspath{{images/}}
\usepackage{fancyhdr}
\usepackage{enumerate}
\usepackage[letterpaper, margin=1in]{geometry}
\usepackage[T1]{fontenc}
\usepackage{newpxtext,newpxmath}
\usepackage{siunitx}
\usepackage{hyperref}
\usepackage{bm}
\usepackage{braket}
\usepackage[version=4]{mhchem}
\usepackage[yyyymmdd,hhmmss]{datetime}  % Insert "Compiled on ..."
\usepackage[symbol]{footmisc}  % footnote symbol

\sisetup{inter-unit-product = \ensuremath{{}\cdot{}}}
\DeclareTextCommand{\nobreakspace}{T1}{\leavevmode\nobreak\ }

\title{Hydrogen, Helium, and \ce{H2} within DFT}		% Title
\author{}								% Author
\date{}											% Date

\pagestyle{fancy}
\fancyhf{}
\rhead{\theauthor}
\lhead{\thetitle}
\lfoot{Last compiled on {\ddmmyyyydate\today} at \currenttime}
\rfoot{\thepage}


\newcommand*{\rmn}[1]{\romannumeral#1}
\newcommand*{\RMN}[1]{\uppercase\expandafter{\romannumeral#1}}
\newcommand*{\tran}{^{\mkern-1.5mu\mathsf{T}}}
\renewcommand{\thefootnote}{\fnsymbol{footnote}}

\begin{document}

%%%%%%%%%%%%%%%%%%%%%%%%%%%%%%%%%%%%%%%%%%%%%%%%%%%%%%%%%%%%%%%%%%%%%%%%%%%%%%%%%%%%%%%%%

\section{Problem \RMN{1}}
\begin{enumerate}[(a)]
	\item First only consider $\varepsilon_{\text{xc}}=\varepsilon_{\text{x}}
	      =-\frac{3}{4} \big( \frac{3}{\pi} \big)^{\frac{1}{3}} n^{ \frac{1}{3} }$,
	      thus $v_{\text{xc}} = - \big( \frac{3}{\pi} \big)^{\frac{1}{3}}%
	      n^{\frac{1}{3}}$.
	      Then Kohn--Sham equation is simply
	      \begin{equation}
	      	\bigg[- \frac{1}{2} \nabla^2 - \frac{1}{r} +
	      	\int \frac{n(\bm{r}')}{\lvert \bm{r} - \bm{r}' \rvert} \, d\bm{r}'
	      	- \Big( \frac{3}{\pi}\Big)^{\frac{1}{3}} n^{\frac{1}{3}}(\bm{r})
	      	\bigg]
	      	\psi(\bm{r}) = \varepsilon \psi(\bm{r}),
	      \end{equation}
	      where $\psi(\bm{r}) = \sum_{i} \phi_i(r)$
	      and for hydrogen atom, $n(\bm{r}) = \psi^\ast(\bm{r}) \psi(\bm{r})$.
	      Left multiplied by $\phi_i^\ast(\bm{r})$ and integrate, we have
	      \begin{equation}
	      	\Braket{\phi_i |
	      		- \frac{1}{2} \nabla^2 - \frac{1}{r} +
	      		\int \frac{n(\bm{r}')}{\lvert \bm{r} - \bm{r}' \rvert} \, d\bm{r}'
	      		- \Big( \frac{3}{\pi}\Big)^{\frac{1}{3}} n^{\frac{1}{3}}(\bm{r})
	      		| \psi } = \varepsilon \braket{\phi_i | \psi}.
	      \end{equation}
	      So the eigenvalue problem is
	      \begin{equation}
	      	\det (H_{ij} + K_{ij} + N_{ij} - E S_{ij}) = 0,
	      \end{equation}
	      where $H_{ij} = \braket{\phi_i | - \frac{1}{2}\nabla^2 - \frac{1}{r} | \phi_j}$, $S_{ij} =
	      \braket{\phi_i | \phi_j}$,
	      \begin{equation}
	      	K_{ij} = 4\pi \int r^2 \, d r \phi_i^\ast ( \bm{r} )
	      	\Big(  \sum_{m, n} c_m c_n
	      	\int  \frac{ \phi_m^\ast  \phi_n^\ast }{ \lvert \bm{r} - \bm{r}' \rvert }\, d \bm{r}' \Big)
	      	\phi_j ( \bm{r} )
	      \end{equation}
	      and
	      \begin{equation}
	      	N_{ij} = - \Big( \frac{3}{\pi}\Big)^{\frac{1}{3}}
	      	\braket{\phi_{i} (\bm{r}) |
	      		n^{\frac{1}{3}}(\bm{r})
	      		| \phi_{j} (\bm{r})}
	      	= -4\pi\Big( \frac{3}{\pi}\Big)^{\frac{1}{3}}
	      	\int r^2 \, d r \phi_{i}^\ast
	      	\Big(\sum_{m, n} c_m c_n  \phi_{m}^\ast  \phi_{n} \Big)^{\frac{1}{3}}
	      	\phi_{j}.
	      \end{equation}
	      Total energy
	      \begin{equation}
	      	E_{\text{total}} = E -
	      	\int \frac{n(\bm{r}')n(\bm{r})}{\lvert \bm{r} - \bm{r}' \rvert} \, d\bm{r}'\, d\bm{r} +
	      	\int n(\bm{r}) \varepsilon_{\text{x}}\, d\bm{r} -
	      	\int n(\bm{r}) v_{\text{x}} \, d\bm{r}.
	      \end{equation}






	      %	      Now consider $\varepsilon_{\text{xc}}=\varepsilon_{\text{x}} + \varepsilon_{\text{c}}$. Then Kohn--Sham equation is
	      %	      \begin{equation}
	      %	      	\bigg[- \frac{1}{2} \nabla^2 - \frac{1}{r} +
	      %	      	\int \frac{n(\bm{r}')}{\lvert \bm{r} - \bm{r}' \rvert} \, d\bm{r}'
	      %	      	- \frac{3}{4}\Big( \frac{3}{\pi}\Big)^{\frac{1}{3}} n^{\frac{1}{3}}(\bm{r})
	      %	      	- v_{\text{c}} \bigg]
	      %	      	\psi(\bm{r}) = \varepsilon \psi(\bm{r}),
	      %	      \end{equation}
	      %
	      %	      \begin{equation}
	      %	      	v_c =
	      %	      	\begin{cases}
	      %	      		\frac{\gamma}{1+
	      %	      		\beta_1 \big( \frac{3}{4\pi n(\bm{r})} \big)^{1/6} +
	      %	      		\beta_2 \big( \frac{3}{4\pi n(\bm{r})} \big)^{\frac{1}{3}}}\thinspace,
	      %	      		  & n(\bm{r}) \leq \frac{3}{4\pi}, \\
	      %	      		\frac{1}{3} A \log\big( \frac{3}{4\pi n(\bm{r})} \big) +
	      %	      		B +\frac{1}{3} \big( \frac{3}{4\pi n(\bm{r})} \big)^{\frac{1}{3}} C\log\big( \frac{3}{4\pi n(\bm{r})} \big)+
	      %	      		\big( \frac{3}{4\pi n(\bm{r})} \big)^{\frac{1}{3}} D,
	      %	      		  & n(\bm{r}) > \frac{3}{4\pi}.
	      %	      	\end{cases}
	      %	      \end{equation}

\end{enumerate}


\end{document}
